\section{Methodology}
\label{sec:methodology}

\emph{Drafting status.} The final methodology will be reconstructed from the
audited data, model, attribution and scoring implementations. Equations will be
checked directly against code before promotion.

The study uses eight CAC and eight GBM tissue sections. The eight GBM sections
belong to four patient groups; sections are therefore paired experimental units
for section-level descriptions but are not eight independent patients. Training
is label-free. Expert masks are used only to evaluate the spatial agreement of
selected-bin ion images.

Each VAE encodes either the central spectrum alone or a contextual input formed
from the central spectrum and neighbouring spectra. The decoder reconstructs
only the centre. Controlled variants include a fixed $3\times3$ neighbour mean,
a $5\times5$ mean, zero context, nonlocal shuffled context and a learned
attention aggregator. These variants answer different questions and are not
treated as interchangeable evidence of spatial causality.

After training, a two-component GMM is fitted to frozen latent means. IG is
computed with the implementation's trapezoidal path integration, not the
right-endpoint approximation used in the previous report draft. Attribution is
aggregated into a ranking of raw spectral bins. Consolidated molecular-peak
views are used for interpretation only; matched-budget mSCF1 evaluates the raw
ranked bins.

For every selected bin, its ion image is compared with an expert mask using
Pearson correlation coefficient (PCC). F1 is evaluated at PCC thresholds 0.3,
0.4, 0.5 and 0.6, and mSCF1 is their unweighted mean. The reproduced S3PL code
uses the nearest position in its ranked correlation list as a threshold-derived
rank cutoff; it is not described as an exactly equivalent literal
$r\geq t$ filter. Independent audits recompute the saved peak scores from raw
spectra and masks.

Computational comparisons report measured native workflow cost under the
specified protocols. Timer boundaries, hardware, repetitions, CPU time, RSS and
sampled GPU memory are stated separately. Because S3PL uses 10 epochs and the
VAE uses 100, these measurements are not interpreted as equal-work evidence of
intrinsic architectural efficiency.

